% Глава: 26-ci
% Статус: задел (кратко), будет дополнено.
\chapter{CI/CD: автоматизация сборки и поставки}
\label{ch:26-ci}

\emph{Раздел будет дополнен в следующих версиях
пособия. Пока --- краткая вводка и ключевые темы.}

\section{Что такое CI/CD}

\textbf{CI} (Continuous Integration) --- автоматический
запуск сборки и тестов на каждый push в
репозиторий. \textbf{CD} (Continuous
Delivery/Deployment) --- автоматическое
подготовка/выкатка релиза. Инструменты: GitHub
Actions, GitLab CI, Jenkins.

Обычно CI работает в облаке (GitHub Actions), но
самохостный runner можно развернуть и на сервере
лаборатории, чтобы сборки проходили в той же среде, что и
ваши ВМ на \fpath{pglab.nsu.ru}.

\section{Типовой пайплайн}

\begin{minted}[bgcolor=gray!6]{bash}
# псевдоструктура CI-конвейера
on: push
jobs:
  build:
    - checkout code
    - build project (make / docker build)
    - run tests
    - (optional) deploy
\end{minted}

\section{Что будет в полном разделе}
\begin{itemize}
  \item Конфигурация GitHub Actions / GitLab CI;
  \item Матрица сборки (несколько версий ОС/зависимостей);
  \item Кэширование зависимостей;
  \item Публикация артефактов и образов;
  \item Деплой в staging и прод.
\end{itemize}

\section{Связь с остальным курсом}

CI опирается на \cmd{make} (\nameref{ch:14-make}),
тестирование
(\nameref{ch:15-testing-optimization}), Docker
(\nameref{ch:19-docker-basics}) и Git
(\nameref{ch:25-git}).

\section{Задачи для самопроверки}
\begin{enumerate}
  \item Что даёт CI команде разработчиков?
  \item Какие этапы обычно есть в пайплайне?
  \item Какую роль играет Docker в CD?
\end{enumerate}
